% !TEX root = ../../main.tex

\section{HarmonyOS 客户端技术}

\subsection{HarmonyOS}

HarmonyOS 是面向多终端场景的新一代智能终端操作系统，其应用开发体系支持一次开发、多端部署，有助于开发者在不同设备形态之间复用应用能力\cite{harmonyos-multi-device-deployment}。
在 HearBridge 系统中，HarmonyOS 移动端是普通用户使用系统的主要入口，需要完成手语资源浏览、动作示范、训练交互、实时识别结果展示、视频选择和句子识别结果展示等功能。

选择 HarmonyOS 作为客户端平台，一方面符合课题“基于 HarmonyOS 系统”的研究目标，另一方面也能够体现国产移动端生态下无障碍辅助训练应用的开发实践价值。
系统通过移动端将手语资源、数字人动作示范、用户练习和识别反馈连接起来，使用户能够在较完整的交互流程中完成手语训练和辅助交流体验。

\subsection{ArkTS}

ArkTS 是 HarmonyOS 应用开发中的主要开发语言，适合与 ArkUI 结合实现声明式、组件化和状态驱动的页面开发。
在 HearBridge 移动端中，页面需要根据用户操作和接口返回结果动态切换不同状态。
例如，在手势训练过程中，页面需要展示当前连接状态、识别标签、置信度和模型版本信息；
在句子视频识别过程中，页面需要根据视频选择、识别中、识别成功、识别失败和语义修正完成等状态更新界面内容。

采用 ArkTS 进行状态驱动开发，有助于将页面展示、用户交互和异步接口结果组织在统一的组件结构中。
对于 HearBridge 这类训练流程较多、状态变化较频繁的移动应用而言，这种开发方式能够提高页面逻辑的清晰度和后续维护效率。

\subsection{ArkUI}

ArkUI 是 HarmonyOS 的声明式 UI 框架，提供文本、图片、列表、按钮、弹窗、滚动容器和布局组件等界面构建能力。
在 HearBridge 移动端中，ArkUI 主要用于构建手语资源列表、训练页面、动作示范区域、识别结果卡片、视频选择区域和句子识别结果展示区域等界面。

系统的移动端界面需要突出训练引导和结果反馈。
手语资源页面帮助用户选择具体练习内容；
训练页面通过数字人资源或动作示范内容辅助用户理解标准动作；
实时识别页面向用户展示当前识别结果；
句子视频识别页面则展示词级识别序列、中文映射和语义修正后的自然表达结果。
通过 ArkUI 的组件化界面组织方式，可以较好地支持这些训练交互场景。

\subsection{HTTP 与 WebSocket}

HearBridge 移动端需要同时与 Spring Boot 服务端和 Python 手势识别服务进行通信。
其中，普通业务数据如用户信息、手语资源、模型版本和视频识别结果主要通过 HTTP 接口获取；
实时手势识别和样本采集等连续交互场景则更适合使用 WebSocket 进行双向通信。

在实时识别场景中，移动端持续发送相机图像帧，识别服务返回当前动作标签、置信度、有效帧状态和模型版本等信息，使用户能够获得接近实时的训练反馈。
在句子视频识别场景中，移动端选择并提交本地视频，由服务端和识别服务完成离线分析，再将词序列、中文展示和语义修正结果返回给用户。
HTTP 与 WebSocket 的结合，使系统能够同时支持普通业务操作和实时识别交互。
